\section{Grader taxonomy}
\label{app:graders}

\begin{table}[H]
\centering
\footnotesize
\setlength{\tabcolsep}{1.5pt}
\begin{tabular}{@{}lcll@{}}
\toprule
\textbf{Suite} & \textbf{Loc.} & \textbf{Robot} & \textbf{Tasks} \\
\midrule
Cross-skill       & 5.1 & SO-101       & 4 reach, 10 pick \\
Multi-task         & 5.2 & SO-101       & 6 dirs + 4 held-out \\
Hard suite        & 5.3 & SO-101       & 12 phys + 1 cov \\
Simple suite      & 5.4 & Piper, UR5e  & 5 short tasks \\
Piper Pick        & 5.4 & Piper        & 3 reach + 3 pick \\
OpenVLA reach     & \ref{app:openvla-sweep} & SO-101 & 5 reach (RGB) \\
Locomotion v2     & 5.4 & Go2, ANYmal-C & 10 loco tasks \\
\bottomrule
\end{tabular}
\caption{Downstream task suites (\S\ref{sec:capability}). Reach suites
run 10 trials per offset; the hard suite is 12 physics-graded plus 1
coverage task.}
\label{tab:suites}
\end{table}

\noindent\textbf{Grader taxonomy.} AutoAdapter-Bench implements 42
grader types in \texttt{eval.py}; 31 are exercised by the released
task suites. Twenty-eight of those 31 read \emph{post-step simulator
state}: end-effector pose (\texttt{ee\_at\_midpoint},
\texttt{ee\_waypoint\_trace}), object pose and contact
(\texttt{object\_above\_object}, \texttt{object\_lifted\_by},
\texttt{objects\_swapped}, \texttt{object\_pose\_match\_relative}),
body height (\texttt{body\_height\_ratio}), and base pose/yaw for
locomotion (\texttt{forward\_progress}, \texttt{turn\_to\_heading},
\texttt{drive\_and\_turn}). The remaining three are deliberately
\emph{behavioral}, and we label them as such rather than fold them
into the physics count: failure-detection
(\texttt{llm\_reports\_failure}, used by \texttt{unreachable\_detect},
where the correct behavior is to \emph{report} an
\texttt{IKUnreachableError} rather than move), tool-liveness
(\texttt{tool\_executed\_without\_crash}, used by
\texttt{gripper\_cycle}), and end-effector coverage
(\texttt{ee\_visited\_n\_positions}, used by
\texttt{visit\_all\_objects}). The first two sit in the simple suite
(\S\ref{sec:cap-cross-robot}); only the third,
\texttt{visit\_all\_objects}, sits in the 13-task hard suite. In the
released code \texttt{ee\_visited\_n\_positions} now verifies that the
end-effector physically came within tolerance of each object's true
world position (falling back to commanded-target coverage only when
the per-call trajectory is unavailable).

\noindent\textbf{Grasp grading.} For grasp tasks,
\texttt{object\_lifted\_by} requires the target to be both lifted
\emph{and} still held within 12\,cm of the end-effector at the final
step, rejecting an object flung by a faulty weld constraint that would
pass a height-only check. The self-contained validation suite includes
a matching \texttt{grasp\_lift} test, so a driver that cannot truly
grasp fails synthesis rather than inflating its Self score.

\section{12-task robustness and statistical tests}
\label{app:robustness}

\noindent\textbf{12-task physics-only robustness.} The single
behavioral task in the hard suite is \texttt{visit\_all\_objects}.
Table~\ref{tab:robustness} reports the SO-101 portability comparison
on the 12 physics-graded tasks alone (dropping that one task), using
the same trials. The model ranking is unchanged: Auto-Adapter beats
zero-shot CaP on six of seven models and beats few-shot CaP on four
(tying Haiku), exactly as on the full 13-task suite
(Table~\ref{tab:portability}). The headline therefore does not depend
on the behavioral task or on how it is graded.

\begin{table}[H]
\centering
\small
\begin{tabular}{lrrrr}
\toprule
\textbf{Model} & \textbf{A-A} & \textbf{CaP} & \textbf{CaP\textsubscript{fs}} & \textbf{Self} \\
\midrule
Sonnet~4.6    & \textbf{67} & 53 & 50 & -- \\
Opus~4.8      & \textbf{60} & 47 & 45 & 42 \\
Haiku~4.5     & \textbf{58} & 50 & \textbf{58} & 42 \\
Nova~Pro      & \textbf{57} & 53 & 42 & -- \\
DeepSeek~V3.2 & 53          & 10 & \textbf{57} & 22 \\
Ministral~8B  & \textbf{50} & 25 & 38 & -- \\
Qwen3~32B     & 45          & \textbf{70} & 53 & -- \\
\bottomrule
\end{tabular}
\caption{SO-101 portability on the \textbf{12 physics-graded} hard-suite
tasks (\texttt{visit\_all\_objects} excluded), \%, $N{=}5$ trial-level.
Bold = best of A-A/CaP/CaP\textsubscript{fs}. Ranking matches the
13-task Table~\ref{tab:portability}: A-A $>$ CaP on 6/7; A-A $>$
CaP\textsubscript{fs} on 4/7 (Haiku tie).}
\label{tab:robustness}
\end{table}

\noindent\textbf{Statistical significance.} \textbf{In brief: the
Auto-Adapter advantage over zero-shot CaP is significant at the
task level, but individual per-model cells are not separately
resolvable at $N{=}5$.} We test Auto-Adapter vs.\
zero-shot CaP with the (model, task) cell as the paired unit (seven
models $\times$ 13 tasks, $n{=}91$ cells; $N{=}5$ trials each), computed
by \texttt{autoadapter\_bench/stats.py}. Auto-Adapter leads by a mean
$+13.0$ points per cell, and a Wilcoxon signed-rank test over the paired
cells is significant ($W{=}138$, $p{<}0.01$; pooled trial rate $0.59$
vs.\ $0.46$). The coarser model-level comparison over the seven
aggregate rates is only suggestive ($+13$ pp mean; sign test $p{=}0.13$,
Wilcoxon $p{=}0.15$; $n{=}7$ is underpowered), so we rest the claim on
the task-level test rather than on the 6/7 win count. Against few-shot
CaP the aggregate effect is \emph{not} significant ($p{=}0.13$), so we
report that comparison descriptively only. Per cell, only the two
largest gaps clear non-overlapping 95\% Wilson intervals (DeepSeek A-A
$0.57\,[0.45,0.68]$ vs.\ CaP $0.09\,[0.04,0.19]$; Ministral A-A
$0.54\,[0.42,0.65]$ vs.\ CaP $0.23\,[0.15,0.35]$); for the four closely matched Anthropic/Nova cells the per-model
intervals overlap. The effect is therefore robust in aggregate but, at
$N{=}5$, individual cell differences below ${\sim}15$ points are not
separately resolvable, so we do not claim per-model significance.

\noindent\textbf{Quadruped agent-model ablation.} The
\S\ref{sec:cap-cross-robot} quadruped numbers use a Sonnet~4.5 agent
on a shared Sonnet~4.5 driver. Re-running the same 10-task suite with
a Sonnet~4.6 agent gives the same qualitative split (ANYmal-C
favors Auto-Adapter, Go2 reverses), consistent with an
embodiment-driven rather than agent-version-driven effect
(Table~\ref{tab:quad-ablation}).

\begin{table}[H]
\centering
\small
\begin{tabular}{llrr}
\toprule
\textbf{Robot} & \textbf{Agent} & \textbf{A-A} & \textbf{CaP} \\
\midrule
ANYmal-C & Sonnet~4.5 & \textbf{68} & 28 \\
ANYmal-C & Sonnet~4.6 & \textbf{62} & 50 \\
Go2      & Sonnet~4.5 & 48 & \textbf{54} \\
Go2      & Sonnet~4.6 & 44 & \textbf{54} \\
\bottomrule
\end{tabular}
\caption{Quadruped locomotion (10-task v2 suite, physics \%, $N{=}5$):
Sonnet~4.5 (main, \S\ref{sec:cap-cross-robot}) vs.\ Sonnet~4.6
(ablation). The ANYmal-favors-A-A / Go2-reverses split holds under
both agent versions.}
\label{tab:quad-ablation}
\end{table}

\subsection{Pretrained-VLA probe: OpenVLA-7B zero-shot on SO-101}
\label{app:openvla-sweep}

As an auxiliary check we ran a pretrained vision-language-action model,
OpenVLA-7B, zero-shot on the five SO-101 reach variants
(50 episodes). This is a cross-embodiment \emph{probe}, not a
like-for-like baseline: OpenVLA reads RGB rather than the privileged
state our agent and the RL/CaP baselines use, SO-101 is not in its
Open-X training mixture, and its camera frame differs from training; we
therefore keep it out of the main comparison. It scores \textbf{0/50}
(mean EE error 23\,cm, deltas in inconsistent directions), consistent
with the per-embodiment finetuning OpenVLA requires
\citep[\S 3]{kim2024openvla}. Because a zero invites the objection that
the action adapter, not the model, failed, we validate the harness and
sweep the two configuration knobs a reviewer would question. \emph{Harness validation:} an oracle that feeds
the ground-truth direction (target $-$ current EE) through the identical
\texttt{move\_cartesian} path reaches all five reach targets (5/5, mean
final error 1.3\,cm) at the per-step control budget used for OpenVLA
(\texttt{duration}{=}0.5\,s); the same path with a shorter
\texttt{duration} (0.1\,s) reaches 0/5. A capable
controller therefore passes on this harness, so a model score of 0 is a
model result. \emph{Robustness:} Table~\ref{tab:openvla-sweep} sweeps
six action scales (0.25--10$\times$) and eight \texttt{unnorm}
de-normalization keys OpenVLA-7B ships (\texttt{bridge\_orig} plus
seven others, the manipulation datasets closest to SO-101). Every cell is $\leq$6.7\% (one chance episode of
15), so the failure is not an artifact of one de-normalization or scale
choice.

\begin{table}[H]
\centering
\small
\setlength{\tabcolsep}{4pt}
\begin{tabular}{llrr}
\toprule
\textbf{unnorm key} & \textbf{scale} & \textbf{success} & \textbf{err (cm)} \\
\midrule
bridge\_orig & 0.25 & 0.0\% & 10.3 \\
bridge\_orig & 0.50 & 0.0\% & 13.1 \\
bridge\_orig & 1.00 & 6.7\% & 15.3 \\
bridge\_orig & 2.00 & 0.0\% & 23.8 \\
bridge\_orig & 5.00 & 0.0\% & 33.7 \\
bridge\_orig & 10.0 & 0.0\% & 47.0 \\
\midrule
fractal20220817\_data & 2.00 & 0.0\% & 58.4 \\
berkeley\_autolab\_ur5 & 2.00 & 0.0\% & 18.7 \\
taco\_play & 2.00 & 0.0\% & 50.2 \\
jaco\_play & 2.00 & 0.0\% & 54.2 \\
viola & 2.00 & 0.0\% & 58.2 \\
toto & 2.00 & 0.0\% & 16.7 \\
austin\_buds & 2.00 & 0.0\% & 57.6 \\
\bottomrule
\end{tabular}
\caption{OpenVLA-7B config-robustness sweep on SO-101 (3 episodes
$\times$ 5 reach variants per cell; success = mean per-variant rate).
Top: action-scale sweep at the canonical key. Bottom: \texttt{unnorm}-key
sweep at the canonical scale. Harness oracle-validated above.}
\label{tab:openvla-sweep}
\end{table}

\subsection{Franka reach per-variant breakdown}
\label{app:franka}

Table~\ref{tab:franka} gives the per-variant Franka reach results
behind \S\ref{sec:scoping-franka} (20/70 overall).

\begin{table}[H]
\centering
\footnotesize
\setlength{\tabcolsep}{4pt}
\renewcommand{\arraystretch}{0.9}
\begin{tabular}{lrr}
\toprule
\textbf{Variant}  & \textbf{Pass} & \textbf{Mean err (cm)} \\
\midrule
reach\_x+5  & 10/10 & 1.76 \\
reach\_z+5  & 10/10 & 0.13 \\
reach\_x-5,~y$\pm$5,~diag &  0/40 & 2.86 \\
reach\_z-5  &  0/10 & 11.07 \\
\textbf{Total} & \textbf{20/70 (28.6\%)} & \textbf{3.49} \\
\bottomrule
\end{tabular}
\caption{Franka reach: 7 variants $\times$ 10 trials;
\texttt{x-5}/\texttt{y$\pm$5}/\texttt{diag} aggregated,
\texttt{z-5} (workspace boundary) separate.}
\label{tab:franka}
\end{table}

\noindent\textbf{Oracle ablation: solver or synthesis?}
\label{app:franka-oracle}
To separate the agent from the control template, a deterministic
scripted oracle issues the \emph{exact} Cartesian target for the same
seven variants under the same methodology and 2\,cm tolerance --- no
LLM anywhere. The oracle reproduces the published failure signature
exactly (2/7; 2--4\,cm on the same four directions; 11.1\,cm on
\texttt{z-5}), so the failures are deterministic template behavior,
not agent errors or synthesis noise (determinism check: the stock and
$5\times$-iteration arms produce bit-identical per-variant errors).
Neither a $5\times$ IK iteration
budget nor a nullspace-regularized DLS variant repairs it (both 2/7),
so a stronger solver configuration alone is not the fix
(Table~\ref{tab:franka-oracle}). A per-variant decomposition shows two
mechanisms: \texttt{z-5} is a true IK failure at the
workspace/joint-limit boundary, while three of the four 2--4\,cm cases
have sub-millimeter IK residuals (\texttt{x-5} is mixed, 1.1\,cm) and
lose accuracy in the execution layer
(in these reaches, which stay near the arm's default configuration,
the wrist joint settles ${\approx}11^\circ$ from its commanded value).
Consistent with an execution-layer component, gravity-compensated
execution of the same stock IK solutions recovers one variant (3/7).
This does not change the reported failure pattern or the ceiling
claim; it narrows the cause from solver choice alone to the broader
prescribed arm-control template.

\begin{table}[H]
\centering
\footnotesize
\setlength{\tabcolsep}{4pt}
\renewcommand{\arraystretch}{0.9}
\begin{tabular}{lrl}
\toprule
\textbf{Arm (deterministic, no LLM)} & \textbf{Pass} & \textbf{Note} \\
\midrule
Stock DLS driver          & 2/7 & same signature as Table~\ref{tab:franka} \\
Stock, $5\times$ iterations & 2/7 & bit-identical errors \\
Nullspace-regularized DLS & 2/7 & fixes \texttt{x-5}, loses \texttt{x+5} \\
Stock + gravity-comp.\ execution & 3/7 & recovers \texttt{diag} \\
\bottomrule
\end{tabular}
\caption{Franka oracle ablation: a scripted oracle replays the exact
targets. The published failures persist with no LLM in the loop and
are not repaired by solver changes alone.}
\label{tab:franka-oracle}
\end{table}

